\section{The nodal gradient reconstruction}
\label{sec:recon}

Rhie--Chow needs a nodal pressure gradient, and the exact higher-order Jacobian needs a nodal
velocity gradient --- the state's for the residual's correction, the direction's for its
derivative. One pass produces all of them
(\code{cvfem\_hex8\_nodal\_grads\_packed\_nc} and its atomic twin in
\code{src/hex8/cvfem\_hex8\_pack\_helpers.hpp}).

\subsection{What it computes}

For each element the reference gradient of the field is contracted from the nodal values and
pushed forward by $\adj\Jmat/\det\Jmat$, giving one element-constant gradient per field
component; that value is then distributed to all eight of the element's nodes with a weight, and
the nodal gradient is the weighted sum over the elements meeting at the node. The route ---
reference gradient, then geometry --- is the same two-stage shape every other gradient in the code
uses, deliberately, so there is one implementation of the mapping.

Three properties shape the kernel:

\begin{itemize}
\item \emph{The element gradient does not depend on the node.} The same $n_c$-vector is
      broadcast to all eight nodes, so the pass is scatter-dominated: at three fields that is
      seventy-two scalar read-modify-writes per element against roughly sixty-six flops of real
      work.
\item \emph{It is a reduction over elements.} Hence the same layout question as the operator:
      pack-local accumulation with a deterministic ghost closure, or atomics.
\item \emph{The field count is a template parameter} (\code{NC}). A runtime trip count in the
      scatter cost measurably, so the single code path is instantiated per field count rather than
      parameterised at runtime.
\end{itemize}

\subsection{One path, two field counts}

The pass is used at $n_c = 1$ (the pressure gradient, in every Rhie--Chow residual) and $n_c = 3$
(the velocity gradient, per matvec in the exact higher-order Jacobian). These were once two
separate implementations, and unifying them removed more than duplication: the multi-field path
could not read strided input, which was the only reason the exact action de-interleaved the
direction's three velocity components into scratch on every matvec. Output is described by a
per-component pointer and stride, input by a stride, so the single kernel reads the interleaved
direction in place.

The unification measured $+10\%$ on the packed layout and $+27\%$ on the standard one at both field
counts. Performance neutrality at $n_c = 1$ was the gate that had to be cleared before anything
else, since that path runs in every Rhie--Chow residual.

\subsection{Rejected optimisations, and why they are recorded}

Three further changes to this pass were built, measured and reverted. They are recorded in the
source with their numbers, because each is the kind of idea that looks obviously right on a static
count and is wrong in fact:

\begin{itemize}
\item \emph{Templating the geometry kind} to hoist the isoparametric branch out of the element
      loop. Lost; and the first diagnosis of why was also wrong, blaming an OpenMP capture that
      turned out to be innocent.
\item \emph{Writing owned rows straight out}, skipping the pack-local buffer for the contiguous
      owned range. The static argument is sound --- it removes a read-and-write pass over the
      owned majority --- and it measured slower.
\item \emph{Lane-blocking the element loop}, which every operator sweep beside it already does.
      Lost at this pass's arithmetic-to-scatter ratio.
\end{itemize}

The reason to keep the numbers rather than the conclusions is that all three arguments remain
superficially convincing, and a reader who re-derives them from a load count will rebuild them.
The same applies to the surface-loop reschedulings of Section~\ref{sec:mk:structure}.

\section{Reading the source from here}
\label{sec:map}

\begin{center}
\begin{tabular}{@{}ll@{}}
\toprule
component & file \\
\midrule
surfaces, geometry, fluxes, packs & \code{src/upwind/cvfem\_hex8\_ns\_upwind\_kernels.hpp} \\
generated higher-order kernels    & \code{src/upwind/cvfem\_hex8\_ns\_upwind\_sympy\_kernels.hpp} \\
their generator                   & \code{python/synthesize\_cvfem\_hex8\_ns\_upwind\_sympy.py} \\
limiters and their derivatives    & \code{src/venkata/cvfem\_venkata\_limiter.hpp} \\
coefficient table, PA tangent, reconstruction & \code{src/hex8/cvfem\_hex8\_pack\_helpers.hpp} \\
packed / standard / coloured sweeps & \code{src/best/cvfem\_hex8\_best\_*.hpp} \\
\bottomrule
\end{tabular}
\end{center}

\noindent
Two conventions that make the source easier to read. Anything uniform over a sweep is a template
parameter, so a function's behaviour is determined by its instantiation and not by its arguments;
and the hand-written scalar kernels are kept as references for the vectorised and generated ones
to be measured against, not as dead code --- they are what the oracles of
Section~\ref{sec:jac} compare against.
